\documentclass[xcolor=svgnames]{beamer}
% ---------------------------------------------------------------
% DeepSeek: Mathematical Foundations & Algorithms  (corrected v2)
% Changes from v1: facts checked against DeepSeek's own papers,
% GRPO slides corrected, one bibliography system (natbib + BibTeX)
% ---------------------------------------------------------------
\usepackage[utf8]{inputenc}
\usepackage{xcolor}
\usepackage{booktabs}
\usepackage{amsmath, amssymb}
\usepackage{graphicx}
\usepackage{tikz}
\usepackage{algorithm}
\usepackage{algpseudocode}
\usepackage{xurl}
\usepackage[numbers,sort&compress]{natbib}
\usetikzlibrary{arrows, positioning}
\usetheme{Madrid}

% COLOURS (same palette as v1)
\definecolor{mqred}{RGB}{166, 25, 46}
\definecolor{mqdeepred}{RGB}{118, 35, 47}
\definecolor{mqgray}{RGB}{55, 58, 54}
\definecolor{mqlightgray}{RGB}{237, 235, 229}
\usecolortheme[named=mqred]{structure}
\setbeamercolor{title in head/foot}{bg=mqlightgray, fg=mqgray}
\setbeamercolor{author in head/foot}{bg=mqdeepred}
\setbeamercolor{page number in head/foot}{bg=mqdeepred, fg=mqlightgray}
\setbeamercolor{background canvas}{bg=white}
\setbeamercolor{block title}{bg=mqdeepred, fg=white}
\setbeamercolor{block body}{bg=mqlightgray!60}
\setbeamertemplate{caption}[numbered]
\setbeamertemplate{bibliography item}[text]
\setbeamerfont{bibliography item}{size=\scriptsize}
\setbeamerfont{bibliography entry author}{size=\scriptsize}
\setbeamerfont{bibliography entry title}{size=\scriptsize}
\setbeamerfont{bibliography entry location}{size=\scriptsize}
\setbeamerfont{bibliography entry note}{size=\scriptsize}

\makeatletter
\setbeamertemplate{footline}{%
  \leavevmode%
  \hbox{%
  \begin{beamercolorbox}[wd=.5\paperwidth,ht=2.25ex,dp=1ex,center]{author in head/foot}%
    \usebeamerfont{author in head/foot}\insertshortauthor~(\insertshortinstitute)
  \end{beamercolorbox}%
  \begin{beamercolorbox}[wd=.4\paperwidth,ht=2.25ex,dp=1ex,center]{title in head/foot}%
    \usebeamerfont{title in head/foot}\insertshorttitle
  \end{beamercolorbox}%
  \begin{beamercolorbox}[wd=.1\paperwidth,ht=2.25ex,dp=1ex,center]{page number in head/foot}%
    \usebeamerfont{page number in head/foot}\insertframenumber{} / \inserttotalframenumber
  \end{beamercolorbox}}%
  \vskip0pt%
}
\makeatother
\beamertemplatenavigationsymbolsempty

\title[DeepSeek: Maths \& Algorithms]{DeepSeek: Mathematical Foundations \& Algorithms}
\subtitle{How Group Relative Policy Optimization (GRPO) teaches LLMs to reason}
\author[Joseph Mwangi Maina]{Joseph Mwangi Maina}
\institute[AIMS Rwanda]{African Institute for Mathematical Sciences (AIMS) Rwanda}
\date{MSc presentation, 2025}
% Logo: add your image file back in Overleaf if you want it on the title slide
% \titlegraphic{\includegraphics[height=1.2cm]{logo.jpeg}}

\begin{document}

\frame{\titlepage}

\begin{frame}{Outline}
  \tableofcontents
\end{frame}

% =================================================================
\section{Introduction}
\begin{frame}{Introduction}
\begin{itemize}
  \item \textbf{DeepSeek} is an AI company founded by \textbf{Liang Wenfeng} on 17 July 2023 in Hangzhou, China, and funded by the hedge fund High-Flyer \cite{sallam2025deepseek}.
  \item \textbf{DeepSeek-R1} (January 2025) is a reasoning model that performs on par with OpenAI's o1 in maths, coding and reasoning \cite{guo2025deepseekr1}.
  \item Open weights under the MIT licence: \textbf{DeepSeek-R1-Zero}, \textbf{DeepSeek-R1}, and six smaller distilled models based on \textbf{Qwen} and \textbf{Llama} \cite{deepseekr1github}.
  \item \textbf{Cost:} R1 is built on DeepSeek-V3, whose final training run used 2.788M H800 GPU hours on a 2,048-GPU cluster, about \textbf{\$5.6M} at \$2 per GPU hour. This excludes earlier research and ablation runs \cite{liu2024deepseek}.
\end{itemize}
\end{frame}

% =================================================================
\section{DeepSeek-V3: the base model}
\begin{frame}{DeepSeek-V3: the base model behind R1}
\begin{itemize}
  \item \textbf{Mixture-of-Experts (MoE):} 671B parameters in total, but only \textbf{37B are active} for each token \cite{liu2024deepseek}.
  \item \textbf{Multi-head Latent Attention (MLA):} compresses keys and values to shrink the memory needed during generation.
  \item \textbf{Multi-token prediction (MTP):} a \emph{training objective} that predicts several future tokens; the extra module can also be reused for \textbf{speculative decoding} at inference.
  \item \textbf{Efficient training:} FP8 mixed-precision training and the DualPipe pipeline-parallel algorithm.
\end{itemize}
\vspace{0.2cm}
\begin{block}{Key idea}
Spend compute only on the experts each token needs, and train in lower precision, so a very large model becomes affordable to train.
\end{block}
\end{frame}

% =================================================================
\section{DeepSeek-R1 versus OpenAI o1}
\begin{frame}{DeepSeek-R1 versus OpenAI o1}
\begin{table}
\centering
\begin{tabular}{lcc}
\toprule
\textbf{Benchmark (metric)} & \textbf{DeepSeek-R1} & \textbf{OpenAI o1-1217} \\
\midrule
AIME 2024 (pass@1) & \textbf{79.8} & 79.2 \\
MATH-500 (pass@1) & \textbf{97.3} & 96.4 \\
Codeforces (percentile) & 96.3 & \textbf{96.6} \\
GPQA Diamond (pass@1) & 71.5 & \textbf{75.7} \\
MMLU (pass@1) & 90.8 & \textbf{91.8} \\
SWE-bench Verified (resolved) & \textbf{49.2} & 48.9 \\
\bottomrule
\end{tabular}
\caption{Reported scores; higher is better \cite{deepseekr1github}.}
\end{table}
R1 matches o1 on maths and coding while being openly released.
\end{frame}

% =================================================================
\section{What DeepSeek is addressing}
\begin{frame}{What DeepSeek is addressing}
\begin{block}{}
\begin{itemize}
  \item \textbf{Reasoning through reinforcement learning:} R1-Zero showed that step-by-step reasoning can emerge from RL with simple rule-based rewards, without supervised fine-tuning first \cite{guo2025deepseekr1}.
  \item \textbf{Training efficiency:} sparse MoE activation and FP8 training cut the compute needed \cite{liu2024deepseek}.
  \item \textbf{Open access:} open weights lower the barrier for developers and organisations, especially outside the largest tech companies \cite{krause2025deepseek}.
\end{itemize}
\end{block}
\end{frame}

% =================================================================
\section{Mathematical foundations}
\begin{frame}{Starting point: Proximal Policy Optimization (PPO)}
For a question $q$ and output tokens $o_t$, PPO maximises \cite{schulman2017ppo}:
\[
J_{\text{PPO}}(\theta) = \mathbb{E}\Big[\min\big(r_t(\theta) A_t,\ \operatorname{clip}(r_t(\theta), 1-\epsilon, 1+\epsilon) A_t\big)\Big],
\]
where the probability ratio between the new and old policy is
\[
r_t(\theta)=\frac{\pi_\theta(o_t \mid q, o_{<t})}{\pi_{\theta_{\text{old}}}(o_t \mid q, o_{<t})}.
\]
\begin{itemize}
  \item $A_t$: advantage, normally estimated with a learned \textbf{value (critic) model}.
  \item The clip with $\epsilon$ stops any single update from moving the policy too far.
\end{itemize}
\begin{alertblock}{The problem for LLMs}
The critic is usually as large as the policy itself, roughly doubling memory and compute \cite{shao2024deepseekmath}.
\end{alertblock}
\end{frame}

\begin{frame}{GRPO: replace the critic with a group baseline}
For each question $q$, sample a \textbf{group} of $G$ answers $\{o_1,\dots,o_G\}$ from the old policy and score them with rewards $r_1,\dots,r_G$ (for example: is the final answer correct, is the format right) \cite{shao2024deepseekmath,guo2025deepseekr1}.

\vspace{0.2cm}
\textbf{Group-relative advantage} (every token of answer $i$ gets the same value):
\[
\hat{A}_{i} = \frac{r_i - \operatorname{mean}(r_1,\dots,r_G)}{\operatorname{std}(r_1,\dots,r_G)}
\]
\begin{itemize}
  \item \textbf{No critic network:} the group mean is the baseline, saving memory.
  \item \textbf{Variance reduction:} answers are judged against their siblings, not on an absolute scale.
  \item Answers better than the group average are pushed up; worse ones are pushed down.
\end{itemize}
\end{frame}

\begin{frame}{The GRPO objective}
\[
J_{\text{GRPO}}(\theta)=\mathbb{E}\Bigg[\frac{1}{G}\sum_{i=1}^{G}\frac{1}{|o_i|}\sum_{t=1}^{|o_i|}
\Big\{\ell_{i,t}(\theta)-\beta\,\mathbb{D}_{\text{KL}}\big[\pi_\theta \,\|\, \pi_{\text{ref}}\big]\Big\}\Bigg],
\]
with the clipped per-token term
\[
\ell_{i,t}(\theta)=\min\big(r_{i,t}(\theta)\hat{A}_{i},\ \operatorname{clip}(r_{i,t}(\theta),1-\epsilon,1+\epsilon)\hat{A}_{i}\big).
\]
\begin{itemize}
  \item The \textbf{KL penalty} (weight $\beta$) keeps the policy close to a reference model and is added to the loss directly, not folded into the reward as in PPO \cite{shao2024deepseekmath}.
  \item KL is estimated per token with an unbiased, always-positive estimator:
  \[
  \mathbb{D}_{\text{KL}} = \frac{\pi_{\text{ref}}}{\pi_\theta} - \log\frac{\pi_{\text{ref}}}{\pi_\theta} - 1 .
  \]
\end{itemize}
\end{frame}

\begin{frame}{Why it works: the policy-gradient view}
The policy gradient weights each token's log-probability by its advantage:
\[
\nabla_\theta J(\theta)=\mathbb{E}\Big[\sum_{t}\nabla_\theta \log \pi_\theta(o_t \mid q,o_{<t})\, \hat{A}_{i}\Big].
\]
\begin{itemize}
  \item Subtracting the group mean is a \textbf{baseline}: it leaves the expected gradient unchanged but lowers its variance.
  \item Dividing by the group standard deviation rescales every question to a similar size of update.
\end{itemize}
\begin{alertblock}{Known caveats (later work)}
The length term $1/|o_i|$ and the standard-deviation scaling introduce biases toward longer wrong answers and toward very easy or very hard questions; ``Dr.~GRPO'' removes them \cite{liu2025drgrpo}.
\end{alertblock}
\end{frame}

% =================================================================
\section{The GRPO algorithm and the R1 pipeline}
\begin{frame}{GRPO: the algorithm}
\begin{algorithm}[H]
\caption{GRPO (simplified, after \cite{shao2024deepseekmath})}
\begin{algorithmic}[1]
\State \textbf{Input:} policy $\pi_\theta$, reference $\pi_{\text{ref}}$, reward function $R$, group size $G$
\For{each training step}
  \State Sample a batch of questions $q$
  \For{each question $q$}
    \State Sample $G$ answers $o_1,\dots,o_G \sim \pi_{\theta_{\text{old}}}(\cdot \mid q)$
    \State Score them: $r_i = R(q, o_i)$
    \State Group-normalise: $\hat{A}_i = (r_i - \text{mean}(r))/\text{std}(r)$
  \EndFor
  \State Update $\theta$ by gradient ascent on $J_{\text{GRPO}}(\theta)$
  \State $\theta_{\text{old}} \gets \theta$
\EndFor
\end{algorithmic}
\end{algorithm}
\end{frame}

\begin{frame}{From R1-Zero to R1}
\begin{itemize}
  \item \textbf{R1-Zero:} GRPO applied directly to the V3 base model with rule-based rewards (answer accuracy and output format). Reasoning emerges, but outputs mix languages and are hard to read \cite{guo2025deepseekr1}.
  \item \textbf{R1} adds a multi-stage pipeline \cite{guo2025deepseekr1}:
  \begin{enumerate}
    \item \emph{Cold start:} a small set of readable long-reasoning examples for supervised fine-tuning.
    \item \emph{Reasoning RL:} GRPO, with an extra \textbf{language-consistency reward}.
    \item \emph{Rejection sampling + SFT:} keep good RL outputs and add general data.
    \item \emph{RL for all scenarios:} a final round for helpfulness and harmlessness.
  \end{enumerate}
  \item \textbf{Distillation:} R1's outputs train smaller Qwen and Llama models \cite{deepseekr1github}.
\end{itemize}
\end{frame}

% =================================================================
\section{Applications}
\begin{frame}{Potential applications}
\begin{itemize}
  \item \textbf{Natural language processing:} generation, translation, sentiment analysis.
  \item \textbf{Scientific research:} step-by-step reasoning for maths and science problems.
  \item \textbf{Healthcare:} supporting analysis of biomedical text and data, with human oversight.
  \item \textbf{Finance:} risk assessment and analysis, where lower inference cost matters \cite{krause2025deepseek}.
\end{itemize}
\end{frame}

% =================================================================
\section{Market reaction and implications}
\begin{frame}{Market reaction and implications}
\begin{block}{Market reaction (27 January 2025)}
\begin{itemize}
  \item Nvidia fell nearly 17\%, losing about \textbf{\$589 billion} in market value, the largest one-day loss in US history; the Nasdaq Composite fell 3.1\% \cite{morales2025nvidia}.
\end{itemize}
\end{block}
\begin{block}{Implications}
\begin{itemize}
  \item Low-cost, open models put pressure on closed providers to cut prices \cite{krause2025deepseek}.
  \item Wider access to strong models, alongside open questions about data security and governance \cite{sallam2025deepseek}.
\end{itemize}
\end{block}
\end{frame}

\section{Outlook}
\begin{frame}{Outlook}
\begin{itemize}
  \item \textbf{Practical AI:} reasoning models moving into healthcare, education and business workflows.
  \item \textbf{Cheaper, open models:} efficiency gains make strong models available beyond large labs.
  \item \textbf{Open research questions:} better reward design, removing GRPO's biases, and making reasoning models calibrated and trustworthy.
\end{itemize}
\end{frame}

% =================================================================
\section*{References}
\begin{frame}[allowframebreaks]{References}
  \scriptsize
  \bibliographystyle{unsrtnat}
  \bibliography{ref}
\end{frame}

\end{document}
